\section{引言}
\label{sec:intro}

安全运营中心收到的告警远多于分析员能调查的数量，而其中大多数最终被证明是良性的~\citep{nodoze}。LLM 智能体带来了希望：它们查询日志、阅读结果，并判断告警的起因~\citep{excytin,clouseau}。不能把遥测数据发给托管模型的组织，可以在自己的网络内用小型开源权重模型运行这样的智能体。这样做把数据留在了内部，却把整个调查交给了一个能力较弱的模型。不过，许多告警属于反复出现的类型：分析员已经知道可能的原因、区分它们的检查，以及以往案例是如何结案的。

一次调查包含两项不同的工作。第一项是\emph{读}：把探针返回的内容（往往是原始日志文本）变成证据，比如“登录失败来自大量家庭宽带地址”。第二项是\emph{决定}：下一步运行哪个探针、证据是否已经足够、接受哪个结论。现有的智能体把两项工作都交给模型~\citep{react}。我们发现，小模型做不好第二项，却是第一项不可或缺的。Qwen2.5-7B 自己选择探针时只解决 8\% 到 13\% 的案例，几乎用光全部预算；Llama-3.1-8B 在全部 3{,}389 次决策中每次都要求再查一个探针，即使它自己的账本显示后验已超过 0.999。然而，当一次日志变更改写了日志行的写法时，固定解析器一行也读不出，同一个 7B 模型却能正确读懂。两半都不能单独工作：直接面对原始日志自行决定时，这个 7B 模型 48 个案例一个也没有解决，尽管只让它给出结果标签时，它能读对同样的文本。

我们的核心想法是，针对分析员已经熟悉的告警类型，据此分工：让模型读，不让模型决定。我们提出 \hesp{}，一个用于本地分诊智能体的控制器。给定告警的候选原因和一组只读探针，\hesp{} 在这些原因上维护一个贝叶斯账本，运行每单位成本期望信息增益最高的探针，在某个原因得到当前证据支持时自行结束调查，并且只接受账本支持的结论。模型提议下一步；当探针返回原始文本时，模型把它解析成账本使用的结果。

这种分工引出三个问题。第一，把决定从模型手中拿走会损失案例吗？不会：在三个场景中的两个，控制器单独就解决了全部案例，放进控制器的任何模型（7B 到 72B）都没有解决更多；强模型唯一的优势是等待确认性观测，而一条停止规则就能复现它。第二，还需要模型吗？只在读取时需要：观测是结构化的时候，控制器不需要模型；日志格式一变，没有模型就每个案例都被升级；而没有控制器，模型面对原始日志一个案例也解决不了。第三，模型读取器安全吗？不安全：攻击者把一致的良性故事写进智能体读取的每一条日志，就能让模型处处采纳它，连跨独立来源的佐证也会被击穿。因此，\hesp{} 记录每条观测由哪个解析器产生，让模型解析出的观测可以确认攻击、但永远不能确认良性结论。另有两条规则保护决定本身：结束守卫拒绝账本不支持的结论，挡住注入日志的指令；良性结论必须由两个独立来源组佐证，挡住单个伪造来源。

我们在三个模拟分诊场景上评估 \hesp{}，其中一个来自公开的 Sigma 检测规则；用五个 7B 到 72B 的开源权重模型，以及三类攻击，包括针对 \hesp{} 本身设计的自适应攻击。这些场景有两个假设，界定了结果的适用范围：候选原因、探针及其结果事先写好，就像分析员为一类已知告警所做的那样；每个原因至少留下一条能把它单独挑出来的观测。在这两个假设下，做决定的是流程而不是模型。当候选原因需要在调查过程中提出，或者证据只有组合起来才成立时，模型是否能参与决定，我们的场景无法回答。

总结起来，本文的贡献如下：
\begin{itemize}
  \item 测量了小型本地模型在告警分诊中失败的位置：它们在选择探针和停止上失败，相近规模的模型失败原因各不相同；而不含任何模型的控制器解决的案例与我们测试的每个模型一样多（\S\ref{sec:decide}）。
  \item 证明了两半都不能单独工作：日志格式改变后，固定解析器让每个案例都被升级，单独决策的 7B 模型一个也没有解决，而同一个模型为控制器读取时解决了其中大部分（\S\ref{sec:read}）。
  \item 提出 \hesp{}，一个让本地 LLM 只读不决定的控制器，用三条规则把攻击者挡在结论之外：读取器信任、结束守卫和跨来源佐证；并在注入指令、伪造来源和针对读取器的自适应攻击下进行了评估（\S\ref{sec:approach}、\S\ref{sec:rawlog}、\S\ref{sec:injection}）。\footnote{代码、预测表、结果文件和回合日志见 \repourl。}
\end{itemize}
